\section{与有限维模的张量积}
\label{sec:03.3.7}

设 $\theta$ 是 $\h^*$ 中的一个 Weyl 群轨道，$\lambda\in\theta$。设 $F$ 是 $\g$ 的有限维表示，$m=\dim F$。设 $(\mu_i;1\le i\le m)$ 是 $F$ 的全部权（按重数计）。由于 $F$ 的权及其重数在 Weyl 群 $W$ 的作用下不变，族 $S(\theta,F)=(\nu_i=\lambda+\mu_i;1\le i\le m)$（$\lambda\in\theta$）只依赖于 $\theta$ 与 $F$。

\begin{lemma}\label{lem:03.7.1}
设 $V\in\mathcal{M}(U_\theta)$，$F$ 是 $\g$ 的有限维表示。则对任意 $\zeta\in\zz(\g)$，元素
$$\prod_{\nu\in S(\theta,F)}\bigl(\zeta-\chi_\nu(\zeta)\bigr)$$
零化 $V\otimes_\C F$。
\end{lemma}

\begin{proof}
设 $\lambda\in\theta$ 反支配。置 $\mathcal{F}=\calO_X\otimes_\C F$，并取 C.2 开头所给的滤过 $(\mathcal{F}_i;1\le i\le m)$。它诱导 $\Delta_\lambda(V)\otimes_{\calO_X}\mathcal{F}$ 的滤过 $(\Delta_\lambda(V)\otimes_{\calO_X}\mathcal{F}_i;1\le i\le m)$，其相应的分次模是 $\Delta_\lambda(V)(\mu_i)$（$1\le i\le m$）的直和。显然，对任意 $\zeta\in\zz(\g)$，$\Delta_\lambda(V)(\mu_i)$ 被 $\zeta-\chi_{\nu_i}(\zeta)$ 零化。由此立即推出，对任意 $\zeta\in\zz(\g)$，$\Delta_\lambda(V)\otimes_{\calO_X}\mathcal{F}$ 被 $\prod_{i=1}^m(\zeta-\chi_{\nu_i}(\zeta))$ 零化。特别地，其整体截面模被这些元素零化。另一方面，由 1.1，
$$\Gamma(X,\Delta_\lambda(V)\otimes_{\calO_X}\mathcal{F})=\Gamma(X,\Delta_\lambda(V)\otimes_\C F)=V\otimes_\C F.$$
\end{proof}

由上述结果特别推出：$V\otimes_\C F$ 有一个由 $U(\g)$-子模构成的有限递增滤过，使所有合成因子都是具有无穷小特征标的模。

设 $\mathcal{M}_{fl}(U(\g))$ 是 $\mathcal{M}(U(\g))$ 中由有限长度 $U(\g)$-模构成的满子范畴。设 $\mathcal{M}_{cc}(U(\g))$ 是 $\mathcal{M}_{fg}(U(\g))$ 的满子范畴，由满足下述条件的 $V\in\mathcal{M}_{fl}(U(\g))$ 组成：对任意有限维 $\g$-模 $F$ 都有 $V\otimes_\C F\in\mathcal{M}_{fl}(U(\g))$。

\begin{remark}\label{rem:03.7.2}
T.~Stafford 的一个例子表明，存在不可约 $U_\theta$-模 $V$，使 $V\otimes_\C F$ 不是 Artin 模。因此 $\mathcal{M}_{cc}(U(\g))$ 严格小于 $\mathcal{M}_{fl}(U(\g))$。
\end{remark}

\begin{lemma}\label{lem:03.7.3}
\begin{enumerate}[label=(\roman*),leftmargin=2.6em,itemsep=0.1em]
\item 范畴 $\mathcal{M}_{cc}(U(\g))$ 是范畴 $\mathcal{M}_{fl}(U(\g))$ 的厚子范畴。
\item 若 $V\in\mathcal{M}_{cc}(U(\g))$ 且 $F$ 是有限维 $\g$-模，则 $V\otimes_\C F\in\mathcal{M}_{cc}(U(\g))$。
\end{enumerate}
\end{lemma}

\begin{proof}
(i) 由 $-\otimes_\C F$ 的正合性立即得到。(ii) 显然。
\end{proof}

\begin{lemma}\label{lem:03.7.4}
设 $\lambda\in\h^*$ 反支配，$\theta=W\cdot\lambda$，$V\in\mathcal{M}(U_\theta)$。则断言
\begin{enumerate}[label=(\roman*),leftmargin=2.6em,itemsep=0.1em]
\item 对任意权 $\nu\in P(\Sigma)$，$U_\theta$-模 $H^p(X,\Delta_\lambda(V)(\nu))$（$p\in\Z_+$）都是有限长度的
\end{enumerate}
蕴含
\begin{enumerate}[label=(\roman*),leftmargin=2.6em,itemsep=0.1em,start=2]
\item 对任意有限维 $\g$-模 $F$，$V\otimes_\C F$ 都是有限长度的 $\g$-模。
\end{enumerate}
若再设 $\lambda$ 正则，则 (i) 与 (ii) 等价。
\end{lemma}

\begin{proof}
沿用 7.1 证明中的记号。由滤过对象的谱序列 [7, III.13.6] 可知，存在一个谱序列，其 $E_1$ 项为
$$H^{p-q}\bigl(X,\Gr_q(\Delta_\lambda(V)\otimes_{\calO_X}\mathcal{F})\bigr)=H^{p-q}\bigl(X,\Delta_\lambda(V)(\nu_q)\bigr),$$
收敛到
$$H^p(X,\Delta_\lambda(V)\otimes_{\calO_X}\mathcal{F})=H^p(X,\Delta_\lambda(V))\otimes_\C F,$$
且后者当 $p=0$ 时等于 $V\otimes_\C F$，当 $p\ne0$ 时为 $0$。由于 $E_1$ 项由有限长度 $\g$-模组成，且所有微分都是 $\g$-模同态，可知 $V\otimes_\C F$ 具有有限长度。

现设 $\lambda$ 正则，且 $V\in\mathcal{M}(U_\theta)$ 满足 (ii)。设 $\mu$ 是支配权，$F$ 是以 $-\mu$ 为最低权的不可约有限维模。则按 C.2.2 与 1.1，
\begin{align*}
\Gamma\bigl(X,\Delta_\lambda(V)(-\mu)\bigr)&=\Gamma\bigl(X,\Delta_\lambda(V)\otimes_{\calO_X}\mathcal{F}\bigr)[\lambda-\mu]\\
&=\bigl(\Gamma(X,\Delta_\lambda(V))\otimes_\C F\bigr)[\lambda-\mu]=(V\otimes_\C F)[\lambda-\mu].
\end{align*}
由 7.3，$V'=(V\otimes_\C F)[\lambda-\mu]$ 也具有性质 (ii)，并且 $\Delta_{\lambda-\mu}(V')=\Delta_\lambda(V)(-\mu)$。
\begin{annote}
底本此处作「由 2」，即指向同一节的 7.2。但 7.2 只是 Stafford 反例的注记，与这里的推理无关；此处所需的恰是 7.3（$\mathcal{M}_{cc}$ 是厚子范畴，且对张量有限维模封闭），由它才能断言 $V'$ 仍满足 (ii)。故按上下文改引 7.3。[来源：底本 p.69，7.4 的证明；与 7.2、7.3 的内容对照]
\end{annote}

由于任意权 $\nu$ 都可写成两个支配权 $\mu'$ 与 $\mu$ 之差，故
$$\Delta_\lambda(V)(\nu)=\Delta_\lambda(V)(\mu'-\mu)=\Delta_{\lambda-\mu}(V')(\mu').$$
因此，只需对支配权 $\nu$ 证明 (i)。

我们对 $p$ 作归纳来完成证明。先设 $p=0$。设 $F$ 是以 $\nu$ 为最高权的不可约有限维 $\g$-模。则对一切 $p$（$0\le p\le\dim X$）有
$$H^p(X,\Delta_\lambda(V)\otimes_{\calO_X}\mathcal{F})=H^p(X,\Delta_\lambda(V))\otimes_\C F,$$
故当 $p>0$ 时左端为零。另一方面，有从 $F_1=\calO(\nu)$ 到 $\mathcal{F}$ 的单射。于是与 $\Delta_\lambda(V)$ 作张量积，得到 $U^\circ$-模的正合列
$$0\longrightarrow\Delta_\lambda(V)(\nu)\longrightarrow\Delta_\lambda(V)\otimes_{\calO_X}\mathcal{F}\longrightarrow K\longrightarrow0.$$
对这个正合列作用 $\Gamma$，可见 $\Gamma(X,\Delta_\lambda(V)(\nu))$ 是张量积 $\Gamma(X,\Delta_\lambda(V))\otimes_\C F=V\otimes_\C F$ 的 $\g$-子模，而后者按假设具有有限长度。这就对 $p=0$ 证明了断言。

现设 $p\ge1$。则由上同调的长正合列可知，$H^p(X,\Delta_\lambda(V)(\nu))$ 是 $H^{p-1}(X,K)$ 的商。另一方面，由 $\mathcal{F}$ 的滤过的定义可知，$K$ 带有自然的 $U^\circ$-模滤过，其相应的分次模 $\Gr K$ 等于 $\bigoplus_\mu\Delta_\lambda(V)(\mu)$，其中和取遍 $F$ 的所有不等于 $\nu$ 的权 $\mu$。由归纳假设，$H^{p-1}(X,\Delta_\lambda(V)(\mu))$ 都是有限长度的 $\g$-模。对 $K$ 的滤过长度作归纳，即知 $H^{p-1}(X,K)$ 是有限长度的 $\g$-模。由此即得断言。
\end{proof}

\begin{proposition}\label{prop:03.7.5}
设 $\lambda\in\h^*$ 反支配，$\theta=W\cdot\lambda$，$V\in\mathcal{M}(U_\theta)$，$F$ 是 $\g$ 的有限维表示。若 $\Delta_\lambda(V)$ 是完整 $\D_\lambda$-模，则 $V\otimes_\C F$ 是有限长度模。
\end{proposition}

\begin{proof}
由 5.3 与 7.4 即得。
\end{proof}

因此，$\mathcal{M}_{cc}(U(\g))$ 包含一切局部化为完整模的 $U_\theta$-模。

\begin{remark}\label{rem:03.7.6}
是否存在 $\mathcal{M}_{cc}(U(\g))$ 中的不可约 $U_\theta$-模 $V$，使它的局部化不完整？
\end{remark}

设 $K(\mathcal{M}_{fl}(U(\g)))$ 是 $\mathcal{M}_{fl}(U(\g))$ 的 Grothendieck 群。记 $\ch$ 为从 $\mathcal{M}_{fl}(U(\g))$ 到 $K(\mathcal{M}_{fl}(U(\g)))$ 的自然特征标映射。

设 $\mathcal{M}_{cc}(U_\theta)$ 是 $\mathcal{M}(U_\theta)$ 的满子范畴，其对象同时属于 $\mathcal{M}_{cc}(U(\g))$。

\begin{theorem}\label{thm:03.7.7}
设 $\lambda\in\h^*$ 正则反支配，$V$ 是 $\mathcal{M}_{cc}(U_\theta)$ 中的不可约模。则存在唯一的函数 $\Phi:P(\Sigma)\to K(\mathcal{M}_{fl}(U(\g)))$，使得
\begin{enumerate}[label=(\roman*),leftmargin=2.6em,itemsep=0.1em]
\item $\Phi(0)=\ch(V)$；
\item $\Phi(\nu)$ 是若干无穷小特征标为 $\chi_{\lambda+\nu}$ 的 $\g$-模之差；
\item 对任意有限维 $\g$-模 $F$，
$$\ch(V\otimes_\C F)=\sum_{\nu}\Phi(\nu),$$
其中求和取遍 $F$ 的全部权 $\nu$（按重数计）。
\end{enumerate}
函数 $\Phi$ 由下述公式给出：对 $\mu\in P(\Sigma)$，
$$\Phi(\mu)=\sum_{p\in\Z_+}(-1)^p\ch\bigl(H^p(X,\Delta_\lambda(V)(\mu))\bigr).$$
\end{theorem}

\begin{proof}
先证明函数
$$\Phi(\mu)=\sum_{p\in\Z_+}(-1)^p\ch\bigl(H^p(X,\Delta_\lambda(V)(\mu))\bigr)$$
具有所要求的三条性质。前两条由定义显然。与 7.4 中一样，由滤过对象的谱序列可得一个谱序列，其 $E_1$ 项为
$$H^{p-q}\bigl(X,\Gr_q(\Delta_\lambda(V)\otimes_{\calO_X}\mathcal{F})\bigr)=H^{p-q}\bigl(X,\Delta_\lambda(V)(\nu_q)\bigr),$$
收敛到
$$H^p(X,\Delta_\lambda(V)\otimes_{\calO_X}\mathcal{F})=H^p(X,\Delta_\lambda(V))\otimes_\C F,$$
且后者当 $p=0$ 时等于 $V\otimes_\C F$，当 $p\ne0$ 时为 $0$。附着于 $E_1$ 项的总复形的 Euler 示性数为
\begin{align*}
\sum_{p,q\in\Z}(-1)^{p+q}\ch\bigl(H^{p-q}(X,\Delta_\lambda(V)(\nu_q))\bigr)
&=\sum_{p,q\in\Z}(-1)^{p-q}\ch\bigl(H^{p-q}(X,\Delta_\lambda(V)(\nu_q))\bigr)\\
&=\sum_{s\in\Z}(-1)^s\ch\bigl(H^s(X,\Delta_\lambda(V)(\nu))\bigr),
\end{align*}
其中最后一个和取遍 $F$ 的全部权 $\nu$（按重数计）。由 Euler 原理，它等于 $E_\infty$ 的总复形的 Euler 示性数，即 $\ch(V\otimes_\C F)$。这就证明了第三条性质。

下面证明唯一性。首先，映射 $F\mapsto\ch(V\otimes_\C F)$ 可延拓到有限维 $\g$-模范畴的 Grothendieck 群上。此外，权与重数的概念可直接搬到这一情形。由 [5, Ch. VI, \S3, no.~4, Prop.~3]，对任意支配权 $\mu$，有限维 $\g$-模范畴的 Grothendieck 群中存在一个对象，其权的集合为 $\{w\mu\mid w\in W\}$，且每个权的重数都是 $1$。因此，$\sum_{w\in W}\Phi(w\mu)$ 由第三条性质唯一确定。

另一方面，$\chi_{\lambda+\mu}=\chi_{\lambda+w\mu}$ 推出存在 $w'\in W$ 使 $w'(\lambda+w\mu)=\lambda+\mu$。这蕴含 $w'\lambda-\lambda=\mu-w'w\mu\in P(\Sigma)$。因此 $w'\in W_\lambda$。由于 $\lambda$ 反支配，$w'\lambda-\lambda$ 是 $\Sigma^+_\lambda$ 中若干根之和。另一方面，$\mu$ 是支配权，而 $\mu-w'w\mu$ 是 $\Sigma^+$ 中若干根之和的相反数。由于 $\lambda$ 正则，这只能是 $w'\lambda=\lambda$ 且 $w'=1$。但这立即推出 $w\mu=\mu$。于是由 (ii)，$\sum_{w\in W}\Phi(w\mu)$ 中的各项对应互不相同的无穷小特征标。这就推出它们被唯一确定。
\end{proof}

映射 $\Phi$ 通常称为附着于 $V$ 的相干族。

\section{单反射的交织函子}
\label{sec:03.3.8}

本节更细致地考察交织函子 $I_{s_\alpha}$（$\alpha\in\Pi$）在不可约 $\D_\lambda$-模上的作用。若 $\check\alpha(\lambda)$ 不是整数，则由 3.22，$I_{s_\alpha}$ 是范畴 $\mathcal{M}_{qc}(\D_\lambda)$ 与 $\mathcal{M}_{qc}(\D_{s_\alpha\lambda})$ 之间的等价。我们想分析的更有趣的情形，是 $\check\alpha(\lambda)$ 为整数。

先从一条简单的几何预备结果开始。

\begin{lemma}\label{lem:03.8.1}
簇 $Z_w$（$w\in W$）都仿射嵌入\index{fangsheqianru@仿射嵌入|hyperpage} $X\times X$。
\end{lemma}

\begin{proof}
簇 $X\times X$ 是 $\g\times\g$ 的旗簇。由 3.1(ii)，$Z_w$（$w\in W$）是对角作用下的 $\operatorname{Int}(\g)$-轨道，故由 H.1.1 它们仿射嵌入。
\end{proof}

设 $\alpha\in\Pi$。记 $X_\alpha$ 为 $\alpha$ 型抛物子代数构成的簇，$p_\alpha$ 为 $X$ 到 $X_\alpha$ 的自然投影。设 $Y_\alpha=X\times_{X_\alpha}X$ 是 $X$ 关于态射 $p_\alpha$ 与自身的纤维积。记 $q_1$ 与 $q_2$ 分别为 $Y_\alpha$ 到第一、第二因子的相应投影。则下图交换：
\begin{center}
{\small\begin{tikzcd}[column sep=small,row sep=small]
Y_\alpha \arrow[r,"q_2"] \arrow[d,"q_1"'] & X \arrow[d,"p_\alpha"]\\
X \arrow[r,"p_\alpha"'] & X_\alpha
\end{tikzcd}}
\end{center}

此外，$Y_\alpha$ 自然嵌入 $X\times X$。这一嵌入把 $Y_\alpha$ 等同于 $X\times X$ 中的闭子簇 $Z_1\cup Z_{s_\alpha}$。在这一等同下，$Z_1$ 是 $Y_\alpha$ 的闭子簇，$Z_{s_\alpha}$ 是 $Y_\alpha$ 的稠密开子簇。此外，由 8.1，$Z_{s_\alpha}$ 仿射嵌入 $Y_\alpha$。

固定基点 $y\in X_\alpha$，记 $P_{\alpha,y}$ 为 $y$ 在 $G$ 中的稳定化子。如我们在……中所讨论的，$X_\alpha$ 上的 $G$-齐性扭微分算子层由 $P_{\alpha,y}$ 的李代数 $\p_{\alpha,y}$ 上的 $P_{\alpha,y}$-不变线性形式参数化。由于 $P_{\alpha,y}$ 连通，线性形式 $\mu\in\p_{\alpha,y}^*$ 是 $P_{\alpha,y}$-不变的，当且仅当它是 $\p_{\alpha,y}$-不变的。因此 $\mu$ 是 $P_{\alpha,y}$-不变的，当且仅当它在 $\p_{\alpha,y}$ 的换位子代数 $[\p_{\alpha,y},\p_{\alpha,y}]$ 上为零。设 $\bb$ 是 $\g$ 中含于 $\p_{\alpha,y}$ 的一个 Borel 子代数。则 $[\p_{\alpha,y},\p_{\alpha,y}]$ 包含 $\bb$ 的幂零根 $\n$，并且有自然映射 $\h\to\p_{\alpha,y}/[\p_{\alpha,y},\p_{\alpha,y}]$。该映射是满射，其核由 $\alpha$ 的对偶根 $\check\alpha$ 张成。因此，$X_\alpha$ 上的 $G$-齐性扭微分算子层由满足 $\check\alpha(\mu)=0$ 的线性形式 $\mu\in\h^*$ 参数化。

此外，对任意满足 $\check\alpha(\mu)=0$ 的 $\mu\in\h^*$，扭微分算子层 $(\D_{X_\alpha,\mu})_{p_\alpha}$ 是 $X$ 上的 $G$-齐性扭微分算子层，并且
$$(\D_{X_\alpha,\mu})_{p_\alpha}=\D_{X,\mu}=\D_{\mu-\rho}.$$

对任意 $\lambda\in\h^*$，$(\D_\lambda)_{q_1}$ 与 $(\D_\lambda)_{q_2}$ 都是 $Y_\alpha$ 上的扭微分算子层。由于 $p_\alpha\circ q_1=p_\alpha\circ q_2$，对任意满足 $\check\alpha(\mu)=0$ 的 $\mu\in\h^*$ 有
$$(\D_{\mu-\rho})_{q_1}=(\D_{\mu-\rho})_{q_2}.$$

设 $\lambda\in\h^*$ 使 $p=-\check\alpha(\lambda)$ 为整数。则可置 $\mu=\lambda+p\rho$。此时 $\check\alpha(\mu)=\check\alpha(\lambda)+p\check\alpha(\rho)=0$，即 $\mu$ 满足上述条件。因此由 ……，
$$(\D_\lambda)_{q_1}=\bigl((\D_{\mu-\rho})\calO((-p+1)\rho)\bigr)_{q_1}
=\bigl((\D_{\mu-\rho})_{q_1}\bigr)_{q_1^*(\calO((-p+1)\rho))};$$
又由于 $\mu=s_\alpha\lambda+p s_\alpha\rho$，同样地，
$$(\D_{s_\alpha\lambda})_{q_1}=\bigl((\D_{\mu-\rho})\calO((-p+1)s_\alpha\rho+\alpha)\bigr)_{q_1}
=\bigl((\D_{\mu-\rho})_{q_1}\bigr)_{q_1^*(\calO((-p+1)s_\alpha\rho+\alpha))}.$$
\begin{annote}
底本这两式的参数与下标有误：第一式右端的内层因子印作 $(\D_{\lambda-\rho})_{q_1}$，第二式的左端印作 $(\D_\lambda)_{q_2}$。按上文 $\mu=\lambda+p\rho$ 的定义，以及随后 8.2 的陈述，二者应分别为 $(\D_{\mu-\rho})_{q_1}$ 与 $(\D_{s_\alpha\lambda})_{q_1}$，已经改正。[来源：底本 p.72，8.2 之前的计算；与 8.2 的陈述对照]
\end{annote}

设 $L$ 是 $Y_\alpha$ 上由
$$L=q_1^*\bigl(\calO((-p+1)s_\alpha\rho+\alpha)\bigr)\otimes_{\calO_{Y_\alpha}}q_2^*\bigl(\calO((-p+1)\rho)\bigr)^{-1}$$
给出的可逆 $\calO_{Y_\alpha}$-模。则按上述计算，得到下述结果。

\begin{lemma}\label{lem:03.8.2}
设 $\lambda\in\h^*$ 使 $p=-\check\alpha(\lambda)$ 为整数。则
$$(\D_{s_\alpha\lambda})_{q_1}=\bigl((\D_\lambda)_{q_2}\bigr)_L.$$
\end{lemma}

特别地，有定义良好的函子
$$U^j:V\longmapsto R^jq_{1+}\bigl(q_2^+(V)\otimes_{\calO_{Y_\alpha}}L\bigr),$$
它们从 $\mathcal{M}_{qc}(\D_\lambda)$ 映入 $\mathcal{M}_{qc}(\D_{s_\alpha\lambda})$。由于 $q_1$ 的纤维是一维的，故当 $j\ne-1,0,1$ 时 $U^j=0$。

现在分析函子 $U^j$ 与交织函子 $I_{s_\alpha}$ 之间的联系。记 $i_1$ 为 $Z_1$ 到 $Y_\alpha$ 的自然包含，$i_\alpha$ 为 $Z_{s_\alpha}$ 到 $Y_\alpha$ 的自然包含。由于 $i_1$ 是闭浸入\index{bijinru@闭浸入|hyperpage}而 $i_\alpha$ 是开仿射浸入，对任意 $W^\bullet\in D^b((\D_{s_\alpha\lambda})_{q_1})$，在 $D^b((\D_{s_\alpha\lambda})_{q_1})$ 中有特异三角
$$\Gamma_{[Z_1]}(W^\bullet)\longrightarrow W^\bullet\longrightarrow i_{\alpha*}\bigl(W^\bullet|_{Z_{s_\alpha}}\bigr)\longrightarrow\Gamma_{[Z_1]}(W^\bullet)[1].$$
由此得到 $D^b(\D_{s_\alpha\lambda})$ 中的特异三角
$$Rq_{1+}\bigl(\Gamma_{[Z_1]}(W^\bullet)\bigr)\longrightarrow Rq_{1+}(W^\bullet)\longrightarrow Rq_{1+}\bigl(i_{\alpha*}(W^\bullet|_{Z_{s_\alpha}})\bigr)\longrightarrow Rq_{1+}\bigl(\Gamma_{[Z_1]}(W^\bullet)[1]\bigr).$$
此外，$q_1\circ i_\alpha=p_1$，故
$$Rq_{1+}\bigl(i_{\alpha*}(W^\bullet|_{Z_{s_\alpha}})\bigr)=Rp_{1+}\bigl(W^\bullet|_{Z_{s_\alpha}}\bigr).$$

现设对某个 $V\in\mathcal{M}_{qc}(\D_\lambda)$ 有 $W^\bullet=D\bigl(q_2^+(V)\otimes_{\calO_{Y_\alpha}}L\bigr)$。则按 3.3(i)，
$$L|_{Z_{s_\alpha}}=p_1^*\bigl(\calO((-p+1)s_\alpha\rho+\alpha)\bigr)\otimes_{\calO_{Z_{s_\alpha}}}p_2^*\bigl(\calO((-p+1)\rho)\bigr)^{-1}=p_1^*(\calO(\alpha))=T_{s_\alpha}.$$
由此可见，此时 $W^\bullet|_{Z_{s_\alpha}}=D\bigl(p_2^+(V)\otimes_{\calO_{Z_{s_\alpha}}}T_{s_\alpha}\bigr)$，于是
$$Rq_{1+}\bigl(i_{\alpha*}(W^\bullet|_{Z_{s_\alpha}})\bigr)=LI_{s_\alpha}(D(V)).$$

由 [1, 7.12] 知 $R\Gamma_{[Z_1]}(W^\bullet)=Ri_{1+}\bigl(Li_1^+(W^\bullet)[-1]\bigr)$，故
$$Rq_{1+}\bigl(R\Gamma_{[Z_1]}(W^\bullet)\bigr)=R(q_1\circ i_1)_+\bigl(Li_1^+(W^\bullet)[-1]\bigr).$$
这里 $q_1\circ i_1$ 是由第一因子上的投影诱导的、$X\times X$ 中对角 $Z_1$ 与 $X$ 之间的自然同构。若再设 $W^\bullet=D\bigl(q_2^+(V)\otimes_{\calO_{Y_\alpha}}L\bigr)$，则
$$Li_1^+\bigl(D(q_2^+(V)\otimes_{\calO_{Y_\alpha}}L)\bigr)=L(q_2\circ i_1)^+(D(V))\otimes_{\calO_{Z_1}}i_1^*(L).$$
这里 $q_2\circ i_1$ 同样是由第二因子上的投影诱导的、$X\times X$ 中对角 $Z_1$ 与 $X$ 之间的自然同构。若用这个映射把 $Z_1$ 与 $X$ 等同，则 $i_1^*(L)=\calO(p\alpha)$，并且
$$Rq_{1+}\bigl(R\Gamma_{[Z_1]}\bigl(D(q_2^+(V)\otimes_{\calO_{Y_\alpha}}L)\bigr)\bigr)=D\bigl(V(p\alpha)[-1]\bigr).$$

对上述特异三角取上同调的长正合列，最终得到下述结果。

\begin{theorem}\label{thm:03.8.3}
设 $\lambda\in\h^*$ 使 $p=-\check\alpha(\lambda)$ 为整数，$V\in\mathcal{M}_{qc}(\D_\lambda)$。则
\begin{enumerate}[label=(\roman*),leftmargin=2.6em,itemsep=0.1em]
\item $U^{-1}(V)=L^{-1}I_{s_\alpha}(V)$；
\item 有 $\D_{s_\alpha\lambda}$-模的正合列
$$0\longrightarrow U^0(V)\longrightarrow I_{s_\alpha}(V)\longrightarrow V(p\alpha)\longrightarrow U^1(V)\longrightarrow0.$$
\end{enumerate}
\end{theorem}

若 $V$ 不可约，还能说得更多。

\begin{theorem}\label{thm:03.8.4}
设 $\lambda\in\h^*$ 使 $p=-\check\alpha(\lambda)$ 为整数，$V$ 是不可约 $\D_\lambda$-模。则或者
\begin{enumerate}[label=(\roman*),leftmargin=2.6em,itemsep=0.1em]
\item $U^{-1}(V)=U^1(V)=V(p\alpha)$ 且 $U^0(V)=0$，此时 $I_{s_\alpha}(V)=0$ 且 $L^{-1}I_{s_\alpha}(V)=V(p\alpha)$；或者
\item $U^{-1}(V)=U^1(V)=0$，此时 $L^{-1}I_{s_\alpha}(V)=0$，且序列
$$0\longrightarrow U^0(V)\longrightarrow I_{s_\alpha}(V)\longrightarrow V(p\alpha)\longrightarrow0$$
正合。
\end{enumerate}
模 $U^0(V)$ 是 $I_{s_\alpha}(V)$ 的最大的、不等于 $I_{s_\alpha}(V)$ 的拟凝聚 $\D_{s_\alpha\lambda}$-子模。
\end{theorem}

\begin{proof}
由 8.4 的陈述与 4.7 立即可见，(ii) 中的 $U^0(V)$ 可刻画为 $I_{s_\alpha}(V)$ 的最大拟凝聚 $\D_{s_\alpha\lambda}$-子模。

为证明其余断言，先注意：若 $U^{-1}(V)\ne0$，则由 8.3 有 $L^{-1}I_{s_\alpha}(V)\ne0$。于是由 4.6，$I_{s_\alpha}(V)=0$ 且 $L^{-1}I_{s_\alpha}(V)=V(p\alpha)$。再一次应用 8.3，可知 $U^0(V)=0$ 且 $U^{-1}(V)=V(p\alpha)$。

设 $U^{-1}(V)=0$。则由 8.3，$L^{-1}I_{s_\alpha}(V)=0$。于是由 4.6 可见 $I_{s_\alpha}(V)\ne0$。余下只需证明 $U^1(V)=0$。

设 $U^1(V)\ne0$。则由 8.3，$U^1(V)=V(p\alpha)$。我们将证明这导致 $I_{s_\alpha}(V)=0$，从而矛盾。这一论证还会让我们对满足 $U^1(V)\ne0$ 的不可约 $\D_\lambda$-模的结构有所了解。

先给出一条预备结果。设 $\mu\in\h^*$ 是满足 $\check\alpha(\mu)=0$ 的线性形式，$U$ 是 $X_\alpha$ 上的 $\D_{X_\alpha,\mu}$-模。则 $p_\alpha^+(U)$ 是 $X$ 上的 $\D_{\mu-\rho}$-模。

\begin{lemma}\label{lem:03.8.5}
设 $V=p_\alpha^+(U)$，其中 $U\in\mathcal{M}_{qc}(\D_{X_\alpha,\mu})$。则 $I_{s_\alpha}(V)=0$ 且 $L^{-1}I_{s_\alpha}(V)=V(\alpha)$。
\end{lemma}

\begin{proof}
如 3.9 的证明之前的讨论，证明归结为 $\g=\mathfrak{sl}(2,\C)$ 的情形。此时 $V$ 是若干个 $\calO_X$ 的拷贝的直和，断言由 3.24 得出。
\end{proof}

这一结果推出：若 $W$ 是形如 $p_\alpha^+(U)$ 的模的平移，则 $I_{s_\alpha}(W)=0$。另一方面，对图表
\begin{center}
{\small\begin{tikzcd}[column sep=small,row sep=small]
Y_\alpha \arrow[r,"q_2"] \arrow[d,"q_1"'] & X \arrow[d,"p_\alpha"]\\
X \arrow[r,"p_\alpha"'] & X_\alpha
\end{tikzcd}}
\end{center}
应用基变换 ……，立即可见 $U^1(V)$ 是这种模的平移。因此 $V$ 也具有这一性质，于是 $I_{s_\alpha}(V)=0$。这就完成了 8.4 的证明。
\end{proof}

\begin{annote}
底本此处作「这就完成了 5.4 的证明」。然而正在进行的论证属于 8.4 的证明；5.4 是第 5 节中的另一条引理，其证明已在第 5 节内结束。故按上下文改为 8.4。[来源：底本 p.74，8.4 的证明末；与底本 p.62 中 5.4 的证明对照]
\end{annote}

此外，我们还证明了下面的结果。

\begin{proposition}\label{prop:03.8.6}
设 $\lambda\in\h^*$ 使 $p=-\check\alpha(\lambda)$ 为整数，$V$ 是不可约 $\D_\lambda$-模。则下述条件等价：
\begin{enumerate}[label=(\roman*),leftmargin=2.6em,itemsep=0.1em]
\item $I_{s_\alpha}(V)=0$；
\item $V$ 是形如 $p_\alpha^+(U)$ 的模的平移。
\end{enumerate}
\end{proposition}

这一结果与 4.9 结合，给出不可约 $\D_\lambda$-模的整体截面消失的如下简单判别法。

\begin{theorem}\label{thm:03.8.7}
设 $\lambda\in\h^*$ 反支配。设 $S$ 是 $\Pi_\lambda$ 中由与 $\lambda$ 正交的全部根组成的子集。假设 $S\subset\Pi$。设 $V$ 是不可约 $\D_\lambda$-模。则下述条件等价：
\begin{enumerate}[label=(\roman*),leftmargin=2.6em,itemsep=0.1em]
\item $\Gamma(X,V)=0$；
\item 存在 $\alpha\in S$，使 $V$ 是形如 $p_\alpha^+(U)$ 的模的平移。
\end{enumerate}
\end{theorem}
